\section{Results and Analysis}
\label{sec:results}

% ADD RESULTS: Present your own measurements and analysis only.
% ADD UNCERTAINTY CALCULATIONS: Explain propagation, fitting and the meaning
% of reported uncertainties. Give units and consistent significant figures.

\subsection{Laser Wavelength Determination}
\label{sec:wavelength}
% ADD MEASUREMENTS, EQUATIONS AND ANALYSIS: Insert your measurements, analysis
% method, fitted quantities, wavelength estimate and uncertainty calculation.

\subsection{Laser Attenuation and Mean Photon Occupation}
\label{sec:attenuation}
% ADD MEASUREMENTS AND EQUATIONS: Record attenuation settings and measured
% quantities; show your occupation calculation, assumptions and uncertainty.

\subsection{Photon-Counting Statistics and Fano-Factor Analysis}
\label{sec:statistics}
% ADD TABLES, GRAPHS AND ANALYSIS: Present count data, acquisition intervals,
% statistical calculations and your Fano-factor estimate with uncertainty.

\subsection{Low-Flux Double-Slit Interference}
\label{sec:interference}
% ADD GRAPHS AND ANALYSIS: Present your measured interference data with units,
% error bars, any fitted model and the corresponding uncertainty calculations.

% FIGURE EXAMPLE: Save the actual graph in figures/selected/ before enabling.
% \begin{figure}[htbp]
%   \centering
%   \includegraphics[width=0.8\linewidth]{your-graph.pdf}
%   \caption{[Describe the measured quantities, conditions and error bars.]}
%   \label{fig:your-graph}
% \end{figure}
% Refer to Figure~\ref{fig:your-graph} in the results text.

% BOOKTABS TABLE EXAMPLE: Replace placeholders with your own measurements.
% \begin{table}[htbp]
%   \centering
%   \caption{[Describe the measurements and uncertainty convention.]}
%   \label{tab:measurements}
%   \begin{tabular}{lll}
%     \toprule
%     Quantity & Value (unit) & Uncertainty (unit) \\
%     \midrule
%     [Quantity] & [Value] & [Uncertainty] \\
%     \bottomrule
%   \end{tabular}
% \end{table}
% Refer to Table~\ref{tab:measurements} in the results text.
% SI UNIT SYNTAX: \si{\metre} for units; use \SI{<value>}{\metre}
% only after replacing <value> with a measured numerical value.
